    \subsection{Overview of DTR PCS \DTRPcs}

    Checks for compliant behavior of the generator under a grid scenario where there is a
    severe loss of load in the bulk transmission network: 7GW out of 307GW. In this
    case the network side is simulated by means of a large generator (340 GVA) and
    two large loads (300 and 7 GW).

    The grid model and its operational point is as in the following schematic:
    \begin{center}
        \includestandalone[width=0.8\textwidth]{circuit_\DTRPcs.tikz}
    \end{center}
    \begin{center}
        \small \textbf{Note: This schematic is only a reminder of the test setup on the TSO's
        side --- the Producer's side may vary, depending on the user-provided model.}
    \end{center}

    \begin{description}
        \item Initial conditions used at the PDR bus:
        \begin{itemize}
            \item $P = P_{max\_unite}$
            \item $Q = 0$
            \item $U = U_{dim}$
        \end{itemize}
    \end{description}

    \subsubsection{Simulation parameters}

    Solver and parameters used in the simulation:
    \begin{center}
        \begin{tabular}{lc}
            \toprule
           \textbf{Parameter} & \textbf{Value (default)} \\
            \midrule
            \BLOCK{for row in solverPCSI8LoadShedDisturbancePmaxQzero}
            {{row[0]}}         & {{row[1]}}                         \\
            \BLOCK{endfor}
            \bottomrule
        \end{tabular}
    \end{center}

    \subsubsection{Simulation}
    As required by the DTR PCS \DTRPcs, the figures below show the
    following magnitudes:
    \begin{itemize}
        \item Voltage at connection point: modulus of the AC complex voltage at
        the PDR bus.
        \item Reactive power supplied at the connection point.
        \item Active power supplied at the connection point.
        \item Rotor speed.
        \item Magnitude controlled by the AVR.
        \item Network frequency.
    \end{itemize}

    \subsubsection{Simulation results}
    The blue line shows the calculated curve, if a reference curve has been entered it is
    shown in orange.

% For now we won't use floats for figures, to get more precise placement
    \noindent
    \begin{minipage}[t]{0.48\textwidth}
        \centering
        \includegraphics[width=\textwidth]{\Producer_fig_V_PCS_RTE-I8.LoadShedDisturbance.PmaxQzero}
        \begin{minipage}[t]{0.8\textwidth}
            \small Voltage magnitude measured at the PDR bus.
        \end{minipage}
    \end{minipage}
    \hfill
    \begin{minipage}[t]{0.48\textwidth}
        \centering
        \includegraphics[width=\textwidth]{\Producer_fig_Ustator_PCS_RTE-I8.LoadShedDisturbance.PmaxQzero}
        \begin{minipage}[t]{0.8\textwidth}
            \small Stator voltage magnitude. The gray dotted line show
            the AVR setpoint. The red dashed lines show the $\pm 5\%$ margin
            around the AVR setpoint.
        \end{minipage}
    \end{minipage}

    \vspace{0.5cm}

    \noindent
    \begin{minipage}[t]{0.48\textwidth}
        \centering
        \includegraphics[width=\textwidth]{\Producer_fig_P_PCS_RTE-I8.LoadShedDisturbance.PmaxQzero}
        \begin{minipage}[t]{0.8\textwidth}
            \small Real power output P, measured at the PDR bus.
        \end{minipage}
    \end{minipage}
    \hfill
    \begin{minipage}[t]{0.48\textwidth}
        \centering
        \includegraphics[width=\textwidth]{\Producer_fig_Q_PCS_RTE-I8.LoadShedDisturbance.PmaxQzero}
        \begin{minipage}[t]{0.8\textwidth}
            \small Reactive power output Q, measured at the PDR bus.
        \end{minipage}
    \end{minipage}

    \vspace{0.5cm}

    \noindent
    \begin{minipage}[t]{0.48\textwidth}
        \centering
        \includegraphics[width=\textwidth]{\Producer_fig_W_PCS_RTE-I8.LoadShedDisturbance.PmaxQzero}
        \begin{minipage}[t]{0.8\textwidth}
            \small Rotor speed.
        \end{minipage}
    \end{minipage}
    \hfill
    \begin{minipage}[t]{0.48\textwidth}
        \centering
        \includegraphics[width=\textwidth]{\Producer_fig_WRef_PCS_RTE-I8.LoadShedDisturbance.PmaxQzero}
        \begin{minipage}[t]{0.8\textwidth}
            \small Network frequency, in Hz. The dashed horizontal lines mark
            the $\pm$200mHz (green line) and $\pm$250mHz (red line) ranges, respectively (centered on
            the starting value of frequency).
        \end{minipage}
    \end{minipage}
    \\[2\baselineskip]
    Go to  {{ linkPCSI8LoadShedDisturbancePmaxQzero }}


    \subsubsection{Analysis of results}

    \noindent No analysis is needed in PCS \DTRPcs.


    \subsubsection{Compliance checks}

    Compliance checks required by PCS \DTRPcs:

    \begin{minipage}{\linewidth} % because otherwise, the footnote does not show
        \begin{tabular}{lcc}
            \toprule
            \textbf{Check} & \multicolumn{1}{c}{\textbf{value}} & \multicolumn{1}{c}{\textbf{Compliant?}} \\
            \midrule
            \BLOCK{for row in cmPCSI8LoadShedDisturbancePmaxQzero}
            {{row[0]}}     & {{row[1]}}                         & {{row[2]}}                              \\
            \BLOCK{endfor}
            \bottomrule
        \end{tabular}
    \end{minipage}
